% 1.5 空间群
\section{空间群}

在讨论空间群的定义之前，先要说明什么是\textit{布拉维格子}。格子是数学点的集合，其中每个格点的环境在相同取向下完全相同。就晶体而言，构成晶体的原子或分子按规则方式排列（Haüy 1815\textsuperscript{a}–\textsuperscript{d}）；也就是说，晶体内部存在一组数学格点。如果让一位麦克斯韦妖从这个格点组中的任意一点观察晶体，他看到的原子或分子的排布与从其他任何格点看到的完全相同。严格地说，要使每个格点的环境完全相似，数学格子必须是无穷延伸的；真实晶体显然不可能包含这样的无穷格子，但考虑到原子的实际尺寸，它非常接近于无穷格子。我们可以用二维例子来说明格子的概念：图 1.6 中那组按特定方式排列、并被想象为延伸到图外无穷远的点，就构成正方二维布拉维格子 $p$。

\textbf{表 1.5}\quad \textit{点群 432（$O$）的群乘法表}

\noindent\resizebox{\textwidth}{!}{%
\begin{tabular}{|c||c||c||c||c||c||c||c||c||c||c||c||c||c||c||c||c||c||c||c||c||c||c||c||c|}
\hline
 & $E$ & $C_{2x}$ & $C_{2y}$ & $C_{2z}$ & $C_{4x}^{-}$ & $C_{4y}^{-}$ & $C_{4z}^{-}$ & $C_{4x}^{+}$ & $C_{4y}^{+}$ & $C_{4z}^{+}$ & $C_{31}^{-}$ & $C_{32}^{-}$ & $C_{33}^{-}$ & $C_{34}^{-}$ & $C_{31}^{+}$ & $C_{32}^{+}$ & $C_{33}^{+}$ & $C_{34}^{+}$ & $C_{2a}$ & $C_{2b}$ & $C_{2c}$ & $C_{2d}$ & $C_{2e}$ & $C_{2f}$ \\ \hline
$E$ & $E$ & $C_{2x}$ & $C_{2y}$ & $C_{2z}$ & $C_{4x}^{-}$ & $C_{4y}^{-}$ & $C_{4z}^{-}$ & $C_{4x}^{+}$ & $C_{4y}^{+}$ & $C_{4z}^{+}$ & $C_{31}^{-}$ & $C_{32}^{-}$ & $C_{33}^{-}$ & $C_{34}^{-}$ & $C_{31}^{+}$ & $C_{32}^{+}$ & $C_{33}^{+}$ & $C_{34}^{+}$ & $C_{2a}$ & $C_{2b}$ & $C_{2c}$ & $C_{2d}$ & $C_{2e}$ & $C_{2f}$ \\ \hline
$C_{2x}$ & $C_{2x}$ & $E$ & $C_{2z}$ & $C_{2y}$ & $C_{4x}^{+}$ & $C_{2d}$ & $C_{2a}$ & $C_{4x}^{-}$ & $C_{2c}$ & $C_{2b}$ & $C_{32}^{-}$ & $C_{31}^{-}$ & $C_{34}^{-}$ & $C_{33}^{-}$ & $C_{34}^{+}$ & $C_{33}^{+}$ & $C_{32}^{+}$ & $C_{31}^{+}$ & $C_{4z}^{-}$ & $C_{4z}^{+}$ & $C_{4y}^{+}$ & $C_{4y}^{-}$ & $C_{2f}$ & $C_{2e}$ \\ \hline
$C_{2y}$ & $C_{2y}$ & $C_{2z}$ & $E$ & $C_{2x}$ & $C_{2e}$ & $C_{4y}^{+}$ & $C_{2b}$ & $C_{2f}$ & $C_{4y}^{-}$ & $C_{2a}$ & $C_{33}^{-}$ & $C_{34}^{-}$ & $C_{31}^{-}$ & $C_{32}^{-}$ & $C_{32}^{+}$ & $C_{31}^{+}$ & $C_{34}^{+}$ & $C_{33}^{+}$ & $C_{4z}^{+}$ & $C_{4z}^{-}$ & $C_{2d}$ & $C_{2c}$ & $C_{4x}^{-}$ & $C_{4x}^{+}$ \\ \hline
$C_{2z}$ & $C_{2z}$ & $C_{2y}$ & $C_{2x}$ & $E$ & $C_{2f}$ & $C_{2c}$ & $C_{4z}^{+}$ & $C_{2e}$ & $C_{2d}$ & $C_{4z}^{-}$ & $C_{34}^{-}$ & $C_{33}^{-}$ & $C_{32}^{-}$ & $C_{31}^{-}$ & $C_{33}^{+}$ & $C_{34}^{+}$ & $C_{31}^{+}$ & $C_{32}^{+}$ & $C_{2b}$ & $C_{2a}$ & $C_{4y}^{-}$ & $C_{4y}^{+}$ & $C_{4x}^{+}$ & $C_{4x}^{-}$ \\ \hline
$C_{4x}^{-}$ & $C_{4x}^{-}$ & $C_{4x}^{+}$ & $C_{2f}$ & $C_{2e}$ & $C_{2x}$ & $C_{32}^{+}$ & $C_{31}^{-}$ & $E$ & $C_{34}^{+}$ & $C_{33}^{-}$ & $C_{2a}$ & $C_{4z}^{-}$ & $C_{2b}$ & $C_{4z}^{+}$ & $C_{4y}^{+}$ & $C_{2d}$ & $C_{4y}^{-}$ & $C_{2c}$ & $C_{32}^{-}$ & $C_{34}^{-}$ & $C_{31}^{+}$ & $C_{33}^{+}$ & $C_{2y}$ & $C_{2z}$ \\ \hline
$C_{4y}^{-}$ & $C_{4y}^{-}$ & $C_{2c}$ & $C_{4y}^{+}$ & $C_{2d}$ & $C_{31}^{-}$ & $C_{2y}$ & $C_{33}^{+}$ & $C_{34}^{-}$ & $E$ & $C_{32}^{+}$ & $C_{2e}$ & $C_{4x}^{+}$ & $C_{4x}^{-}$ & $C_{2f}$ & $C_{4z}^{+}$ & $C_{2a}$ & $C_{2b}$ & $C_{4z}^{-}$ & $C_{31}^{+}$ & $C_{34}^{+}$ & $C_{2z}$ & $C_{2x}$ & $C_{33}^{-}$ & $C_{32}^{-}$ \\ \hline
$C_{4z}^{-}$ & $C_{4z}^{-}$ & $C_{2b}$ & $C_{2a}$ & $C_{4z}^{+}$ & $C_{34}^{+}$ & $C_{31}^{-}$ & $C_{2z}$ & $C_{33}^{+}$ & $C_{32}^{-}$ & $E$ & $C_{2c}$ & $C_{2d}$ & $C_{4y}^{+}$ & $C_{4y}^{-}$ & $C_{4x}^{+}$ & $C_{4x}^{-}$ & $C_{2e}$ & $C_{2f}$ & $C_{2x}$ & $C_{2y}$ & $C_{34}^{-}$ & $C_{33}^{-}$ & $C_{31}^{+}$ & $C_{32}^{+}$ \\ \hline
$C_{4x}^{+}$ & $C_{4x}^{+}$ & $C_{4x}^{-}$ & $C_{2e}$ & $C_{2f}$ & $E$ & $C_{33}^{+}$ & $C_{32}^{-}$ & $C_{2x}$ & $C_{31}^{+}$ & $C_{34}^{-}$ & $C_{4z}^{-}$ & $C_{2a}$ & $C_{4z}^{+}$ & $C_{2b}$ & $C_{2c}$ & $C_{4y}^{-}$ & $C_{2d}$ & $C_{4y}^{+}$ & $C_{31}^{-}$ & $C_{33}^{-}$ & $C_{34}^{+}$ & $C_{32}^{+}$ & $C_{2z}$ & $C_{2y}$ \\ \hline
$C_{4y}^{+}$ & $C_{4y}^{+}$ & $C_{2d}$ & $C_{4y}^{-}$ & $C_{2c}$ & $C_{33}^{-}$ & $E$ & $C_{34}^{+}$ & $C_{32}^{-}$ & $C_{2y}$ & $C_{31}^{+}$ & $C_{4x}^{-}$ & $C_{2f}$ & $C_{2e}$ & $C_{4x}^{+}$ & $C_{2a}$ & $C_{4z}^{+}$ & $C_{4z}^{-}$ & $C_{2b}$ & $C_{32}^{+}$ & $C_{33}^{+}$ & $C_{2x}$ & $C_{2z}$ & $C_{31}^{-}$ & $C_{34}^{-}$ \\ \hline
$C_{4z}^{+}$ & $C_{4z}^{+}$ & $C_{2a}$ & $C_{2b}$ & $C_{4z}^{-}$ & $C_{32}^{+}$ & $C_{34}^{-}$ & $E$ & $C_{31}^{+}$ & $C_{33}^{-}$ & $C_{2z}$ & $C_{4y}^{-}$ & $C_{4y}^{+}$ & $C_{2d}$ & $C_{2c}$ & $C_{2e}$ & $C_{2f}$ & $C_{4x}^{+}$ & $C_{4x}^{-}$ & $C_{2y}$ & $C_{2x}$ & $C_{31}^{-}$ & $C_{32}^{-}$ & $C_{33}^{+}$ & $C_{34}^{+}$ \\ \hline
$C_{31}^{-}$ & $C_{31}^{-}$ & $C_{34}^{-}$ & $C_{32}^{-}$ & $C_{33}^{-}$ & $C_{2c}$ & $C_{2a}$ & $C_{2e}$ & $C_{4y}^{-}$ & $C_{4z}^{-}$ & $C_{4x}^{-}$ & $C_{31}^{+}$ & $C_{33}^{+}$ & $C_{34}^{+}$ & $C_{32}^{+}$ & $E$ & $C_{2x}$ & $C_{2y}$ & $C_{2z}$ & $C_{4x}^{+}$ & $C_{2f}$ & $C_{4z}^{+}$ & $C_{2b}$ & $C_{4y}^{+}$ & $C_{2d}$ \\ \hline
$C_{32}^{-}$ & $C_{32}^{-}$ & $C_{33}^{-}$ & $C_{31}^{-}$ & $C_{34}^{-}$ & $C_{4y}^{+}$ & $C_{4z}^{-}$ & $C_{2f}$ & $C_{2d}$ & $C_{2a}$ & $C_{4x}^{+}$ & $C_{34}^{+}$ & $C_{32}^{+}$ & $C_{31}^{+}$ & $C_{33}^{+}$ & $C_{2x}$ & $E$ & $C_{2z}$ & $C_{2y}$ & $C_{4x}^{-}$ & $C_{2e}$ & $C_{2b}$ & $C_{4z}^{+}$ & $C_{2c}$ & $C_{4y}^{-}$ \\ \hline
$C_{33}^{-}$ & $C_{33}^{-}$ & $C_{32}^{-}$ & $C_{34}^{-}$ & $C_{31}^{-}$ & $C_{2d}$ & $C_{4z}^{+}$ & $C_{4x}^{-}$ & $C_{4y}^{+}$ & $C_{2b}$ & $C_{2e}$ & $C_{32}^{+}$ & $C_{34}^{+}$ & $C_{33}^{+}$ & $C_{31}^{+}$ & $C_{2y}$ & $C_{2z}$ & $E$ & $C_{2x}$ & $C_{2f}$ & $C_{4x}^{+}$ & $C_{2a}$ & $C_{4z}^{-}$ & $C_{4y}^{-}$ & $C_{2c}$ \\ \hline
$C_{34}^{-}$ & $C_{34}^{-}$ & $C_{31}^{-}$ & $C_{33}^{-}$ & $C_{32}^{-}$ & $C_{4y}^{-}$ & $C_{2b}$ & $C_{4x}^{+}$ & $C_{2c}$ & $C_{4z}^{+}$ & $C_{2f}$ & $C_{33}^{+}$ & $C_{31}^{+}$ & $C_{32}^{+}$ & $C_{34}^{+}$ & $C_{2z}$ & $C_{2y}$ & $C_{2x}$ & $E$ & $C_{2e}$ & $C_{4x}^{-}$ & $C_{4z}^{-}$ & $C_{2a}$ & $C_{2d}$ & $C_{4y}^{+}$ \\ \hline
$C_{31}^{+}$ & $C_{31}^{+}$ & $C_{32}^{+}$ & $C_{33}^{+}$ & $C_{34}^{+}$ & $C_{4z}^{+}$ & $C_{4x}^{+}$ & $C_{4y}^{+}$ & $C_{2a}$ & $C_{2e}$ & $C_{2c}$ & $E$ & $C_{2y}$ & $C_{2z}$ & $C_{2x}$ & $C_{31}^{-}$ & $C_{34}^{-}$ & $C_{32}^{-}$ & $C_{33}^{-}$ & $C_{4y}^{-}$ & $C_{2d}$ & $C_{4x}^{-}$ & $C_{2f}$ & $C_{4z}^{-}$ & $C_{2b}$ \\ \hline
$C_{32}^{+}$ & $C_{32}^{+}$ & $C_{31}^{+}$ & $C_{34}^{+}$ & $C_{33}^{+}$ & $C_{2a}$ & $C_{2f}$ & $C_{4y}^{-}$ & $C_{4z}^{+}$ & $C_{4x}^{-}$ & $C_{2d}$ & $C_{2y}$ & $E$ & $C_{2x}$ & $C_{2z}$ & $C_{33}^{-}$ & $C_{32}^{-}$ & $C_{34}^{-}$ & $C_{31}^{-}$ & $C_{4y}^{+}$ & $C_{2c}$ & $C_{2e}$ & $C_{4x}^{+}$ & $C_{2b}$ & $C_{4z}^{-}$ \\ \hline
$C_{33}^{+}$ & $C_{33}^{+}$ & $C_{34}^{+}$ & $C_{31}^{+}$ & $C_{32}^{+}$ & $C_{4z}^{-}$ & $C_{2e}$ & $C_{2d}$ & $C_{2b}$ & $C_{4x}^{+}$ & $C_{4y}^{-}$ & $C_{2z}$ & $C_{2x}$ & $E$ & $C_{2y}$ & $C_{34}^{-}$ & $C_{31}^{-}$ & $C_{33}^{-}$ & $C_{32}^{-}$ & $C_{2c}$ & $C_{4y}^{+}$ & $C_{2f}$ & $C_{4x}^{-}$ & $C_{4z}^{+}$ & $C_{2a}$ \\ \hline
$C_{34}^{+}$ & $C_{34}^{+}$ & $C_{33}^{+}$ & $C_{32}^{+}$ & $C_{31}^{+}$ & $C_{2b}$ & $C_{4x}^{-}$ & $C_{2c}$ & $C_{4z}^{-}$ & $C_{2f}$ & $C_{4y}^{+}$ & $C_{2x}$ & $C_{2z}$ & $C_{2y}$ & $E$ & $C_{32}^{-}$ & $C_{33}^{-}$ & $C_{31}^{-}$ & $C_{34}^{-}$ & $C_{2d}$ & $C_{4y}^{-}$ & $C_{4x}^{+}$ & $C_{2e}$ & $C_{2a}$ & $C_{4z}^{+}$ \\ \hline
$C_{2a}$ & $C_{2a}$ & $C_{4z}^{+}$ & $C_{4z}^{-}$ & $C_{2b}$ & $C_{31}^{+}$ & $C_{32}^{-}$ & $C_{2y}$ & $C_{32}^{+}$ & $C_{31}^{-}$ & $C_{2x}$ & $C_{4y}^{+}$ & $C_{4y}^{-}$ & $C_{2c}$ & $C_{2d}$ & $C_{4x}^{-}$ & $C_{4x}^{+}$ & $C_{2f}$ & $C_{2e}$ & $E$ & $C_{2z}$ & $C_{33}^{-}$ & $C_{34}^{-}$ & $C_{34}^{+}$ & $C_{33}^{+}$ \\ \hline
$C_{2b}$ & $C_{2b}$ & $C_{4z}^{-}$ & $C_{4z}^{+}$ & $C_{2a}$ & $C_{33}^{+}$ & $C_{33}^{-}$ & $C_{2x}$ & $C_{34}^{+}$ & $C_{34}^{-}$ & $C_{2y}$ & $C_{2d}$ & $C_{2c}$ & $C_{4y}^{-}$ & $C_{4y}^{+}$ & $C_{2f}$ & $C_{2e}$ & $C_{4x}^{-}$ & $C_{4x}^{+}$ & $C_{2z}$ & $E$ & $C_{32}^{-}$ & $C_{31}^{-}$ & $C_{32}^{+}$ & $C_{31}^{+}$ \\ \hline
$C_{2c}$ & $C_{2c}$ & $C_{4y}^{-}$ & $C_{2d}$ & $C_{4y}^{+}$ & $C_{34}^{-}$ & $C_{2x}$ & $C_{31}^{+}$ & $C_{31}^{-}$ & $C_{2z}$ & $C_{34}^{+}$ & $C_{4x}^{+}$ & $C_{2e}$ & $C_{2f}$ & $C_{4x}^{-}$ & $C_{4z}^{-}$ & $C_{2b}$ & $C_{2a}$ & $C_{4z}^{+}$ & $C_{33}^{+}$ & $C_{32}^{+}$ & $E$ & $C_{2y}$ & $C_{32}^{-}$ & $C_{33}^{-}$ \\ \hline
$C_{2d}$ & $C_{2d}$ & $C_{4y}^{+}$ & $C_{2c}$ & $C_{4y}^{-}$ & $C_{32}^{-}$ & $C_{2z}$ & $C_{32}^{+}$ & $C_{33}^{-}$ & $C_{2x}$ & $C_{33}^{+}$ & $C_{2f}$ & $C_{4x}^{-}$ & $C_{4x}^{+}$ & $C_{2e}$ & $C_{2b}$ & $C_{4z}^{-}$ & $C_{4z}^{+}$ & $C_{2a}$ & $C_{34}^{+}$ & $C_{31}^{+}$ & $C_{2y}$ & $E$ & $C_{34}^{-}$ & $C_{31}^{-}$ \\ \hline
$C_{2e}$ & $C_{2e}$ & $C_{2f}$ & $C_{4x}^{+}$ & $C_{4x}^{-}$ & $C_{2z}$ & $C_{31}^{+}$ & $C_{33}^{-}$ & $C_{2y}$ & $C_{33}^{+}$ & $C_{31}^{-}$ & $C_{4z}^{+}$ & $C_{2b}$ & $C_{4z}^{-}$ & $C_{2a}$ & $C_{4y}^{-}$ & $C_{2c}$ & $C_{4y}^{+}$ & $C_{2d}$ & $C_{34}^{-}$ & $C_{32}^{-}$ & $C_{32}^{+}$ & $C_{34}^{+}$ & $E$ & $C_{2x}$ \\ \hline
$C_{2f}$ & $C_{2f}$ & $C_{2e}$ & $C_{4x}^{-}$ & $C_{4x}^{+}$ & $C_{2y}$ & $C_{34}^{+}$ & $C_{34}^{-}$ & $C_{2z}$ & $C_{32}^{+}$ & $C_{32}^{-}$ & $C_{2b}$ & $C_{4z}^{+}$ & $C_{2a}$ & $C_{4z}^{-}$ & $C_{2d}$ & $C_{4y}^{+}$ & $C_{2c}$ & $C_{4y}^{-}$ & $C_{33}^{-}$ & $C_{31}^{-}$ & $C_{33}^{+}$ & $C_{31}^{+}$ & $C_{2x}$ & $E$ \\ \hline
\end{tabular}}

\parbox{0.92\textwidth}{
\textit{表 1.5 的注.} 列标题自左至右依次为：$E$；$C_{2x}, C_{2y}, C_{2z}$；$C_{4x}^{-}, C_{4y}^{-}, C_{4z}^{-}, C_{4x}^{+}, C_{4y}^{+}, C_{4z}^{+}$；$C_{31}^{-}, C_{32}^{-}, C_{33}^{-}, C_{34}^{-}, C_{31}^{+}, C_{32}^{+}, C_{33}^{+}, C_{34}^{+}$；$C_{2a}, C_{2b}, C_{2c}, C_{2d}, C_{2e}, C_{2f}$。行标题相同，顺序一致。求乘积 $LM$ 时，取行首元素为 $L$ 的一行与列标题为 $M$ 的一列的交叉格。——中译者
}

\textbf{表 1.6}\quad \textit{点群 622（$D_{6}$）的群乘法表}

\noindent\resizebox{\textwidth}{!}{%
\begin{tabular}{|c||c||c||c||c||c||c||c||c||c||c||c||c|}
\hline
 & $E$ & $C_{6}^{+}$ & $C_{3}^{+}$ & $C_{2}$ & $C_{3}^{-}$ & $C_{6}^{-}$ & $C'_{21}$ & $C'_{22}$ & $C'_{23}$ & $C''_{21}$ & $C''_{22}$ & $C''_{23}$ \\ \hline
$E$ & $E$ & $C_{6}^{+}$ & $C_{3}^{+}$ & $C_{2}$ & $C_{3}^{-}$ & $C_{6}^{-}$ & $C'_{21}$ & $C'_{22}$ & $C'_{23}$ & $C''_{21}$ & $C''_{22}$ & $C''_{23}$ \\ \hline
$C_{6}^{+}$ & $C_{6}^{+}$ & $C_{3}^{+}$ & $C_{2}$ & $C_{3}^{-}$ & $C_{6}^{-}$ & $E$ & $C''_{22}$ & $C''_{23}$ & $C''_{21}$ & $C'_{22}$ & $C'_{23}$ & $C'_{21}$ \\ \hline
$C_{3}^{+}$ & $C_{3}^{+}$ & $C_{2}$ & $C_{3}^{-}$ & $C_{6}^{-}$ & $E$ & $C_{6}^{+}$ & $C'_{23}$ & $C'_{21}$ & $C'_{22}$ & $C''_{23}$ & $C''_{21}$ & $C''_{22}$ \\ \hline
$C_{2}$ & $C_{2}$ & $C_{3}^{-}$ & $C_{6}^{-}$ & $E$ & $C_{6}^{+}$ & $C_{3}^{+}$ & $C''_{21}$ & $C''_{22}$ & $C''_{23}$ & $C'_{21}$ & $C'_{22}$ & $C'_{23}$ \\ \hline
$C_{3}^{-}$ & $C_{3}^{-}$ & $C_{6}^{-}$ & $E$ & $C_{6}^{+}$ & $C_{3}^{+}$ & $C_{2}$ & $C'_{22}$ & $C'_{23}$ & $C'_{21}$ & $C''_{22}$ & $C''_{23}$ & $C''_{21}$ \\ \hline
$C_{6}^{-}$ & $C_{6}^{-}$ & $E$ & $C_{6}^{+}$ & $C_{3}^{+}$ & $C_{2}$ & $C_{3}^{-}$ & $C''_{23}$ & $C''_{21}$ & $C''_{22}$ & $C'_{23}$ & $C'_{21}$ & $C'_{22}$ \\ \hline
$C'_{21}$ & $C'_{21}$ & $C''_{23}$ & $C'_{22}$ & $C''_{21}$ & $C'_{23}$ & $C''_{22}$ & $E$ & $C_{3}^{+}$ & $C_{3}^{-}$ & $C_{2}$ & $C_{6}^{-}$ & $C_{6}^{+}$ \\ \hline
$C'_{22}$ & $C'_{22}$ & $C''_{21}$ & $C'_{23}$ & $C''_{22}$ & $C'_{21}$ & $C''_{23}$ & $C_{3}^{-}$ & $E$ & $C_{3}^{+}$ & $C_{6}^{+}$ & $C_{2}$ & $C_{6}^{-}$ \\ \hline
$C'_{23}$ & $C'_{23}$ & $C''_{22}$ & $C'_{21}$ & $C''_{23}$ & $C'_{22}$ & $C''_{21}$ & $C_{3}^{+}$ & $C_{3}^{-}$ & $E$ & $C_{6}^{-}$ & $C_{6}^{+}$ & $C_{2}$ \\ \hline
$C''_{21}$ & $C''_{21}$ & $C'_{23}$ & $C''_{22}$ & $C'_{21}$ & $C''_{23}$ & $C'_{22}$ & $C_{2}$ & $C_{6}^{-}$ & $C_{6}^{+}$ & $E$ & $C_{3}^{+}$ & $C_{3}^{-}$ \\ \hline
$C''_{22}$ & $C''_{22}$ & $C'_{21}$ & $C''_{23}$ & $C'_{22}$ & $C''_{21}$ & $C'_{23}$ & $C_{6}^{+}$ & $C_{2}$ & $C_{6}^{-}$ & $C_{3}^{-}$ & $E$ & $C_{3}^{+}$ \\ \hline
$C''_{23}$ & $C''_{23}$ & $C'_{22}$ & $C''_{21}$ & $C'_{23}$ & $C''_{22}$ & $C'_{21}$ & $C_{6}^{-}$ & $C_{6}^{+}$ & $C_{2}$ & $C_{3}^{+}$ & $C_{3}^{-}$ & $E$ \\ \hline
\end{tabular}}

\parbox{0.92\textwidth}{
\textit{表 1.6 的注.} 列标题自左至右依次为：$E$；$C_{6}^{+}, C_{3}^{+}, C_{2}, C_{3}^{-}, C_{6}^{-}$；$C'_{21}, C'_{22}, C'_{23}, C''_{21}, C''_{22}, C''_{23}$。行标题相同，顺序一致。求乘积 $LM$ 时，取行首元素为 $L$ 的一行与列标题为 $M$ 的一列的交叉格。——中译者
}

这些点与图 1.6 中所画的完全相似。严格地说，要使每个格点的环境完全相似，数学格子必须是无穷延伸的；真实晶体显然不可能包含这样的无穷格子，但考虑到原子的实际尺寸，它非常接近于无穷格子。我们可以用二维例子来说明格子的概念：若图 1.6 中排列在正方形角点上的那组点被想象为延伸到图外无穷远，则这些点构成简单正方二维布拉维格子 $p$。二维布拉维格子共有五种（Alexander 1929，Alexander and Herrmann 1929，Belov 1959\textsuperscript{b}），它们在《X 射线晶体学国际表》第一卷中有描述并绘出（Henry and Lonsdale 1965）。在三维欧氏空间 $E(3)$ 中，本质上不同的格子有 14 种，称为（三维）\textit{布拉维格子}（Bravais 1849\textsuperscript{a}, \textsuperscript{b}, 1850）。或许应当强调：真实晶体中并非每个原子中心都一定有布拉维格点，也并不要求布拉维格点上一定有原子。

若取某个特定格点为原点 $O$，则其他任何格点的位置向量 $\mathbf{t}$ 可写为

\begin{equation*}
\mathbf{t} = n_{1}\mathbf{t}_{1} + n_{2}\mathbf{t}_{2} + n_{3}\mathbf{t}_{3}
\tag{1.5.1}
\end{equation*}

其中 $n_{1}$、$n_{2}$、$n_{3}$ 是整数，$\mathbf{t}_{1}$、$\mathbf{t}_{2}$、$\mathbf{t}_{3}$ 是连接 $O$ 与其三个（非共面的）最近邻的基本平移。由式 (1.5.1) 定义的平移操作 $\mathbf{t}$ 构成一个无穷群，称为布拉维格子的\textit{平移群}（translation group）。除平移群的对称性之外，布拉维格子还对作用在格点上（有时至少也作用在其他点上）的某些点群操作保持不变。作用在任一格点上的点群操作构成一个点群 $\mathbf{P}$，它实际上是七大晶系之一的全对称点群，即 $\mathbf{P}$ 是该晶系中对称操作数目最多的点群。在图 1.6 所示的二维布拉维格子的例子里，格子的对称性为点群 $4mm$（$C_{4v}$）；若把点画在纸的两面，则为 $4/mmm$（$D_{4h}$）。七个晶系——三斜、单斜、正交、四方、三方、六方、立方——的全对称点群 $\mathbf{P}$ 分别是 $\bar{1}$（$C_{i}$）、$2/m$（$C_{2h}$）、$mmm$（$D_{2h}$）、$4/mmm$（$D_{4h}$）、$\bar{3}m$（$D_{3d}$）、$6/mmm$（$D_{6h}$）与 $m3m$（$O_{h}$）。$\mathbf{P}$ 是相应晶系的全对称点群，源于以下事实：(i) $\mathbf{P}$ 必须包含空间反演操作，因为若 $\mathbf{t}$ 是布拉维格子的向量，则 $-\mathbf{t}$ 也是；(ii) 与 3、4、6 阶轴一起，$\mathbf{P}$ 还包含过这些轴的反射面。因此可以把 14 个布拉维格子归入七个晶系之一，这些布拉维格子绘于图 1.7。根本上说，正是由于晶体固体的点群对称性必须与该固体布拉维格子的对称性相容，才只有 32 个可能的晶体学点群：32 个晶体学点群中的每一个，要么是某个布拉维格子的点群（即七个晶系之一的全对称点群），要么是这七个点群之一的子群。

\begin{figure}[H]
\centering
\includegraphics[width=0.62\textwidth]{fig1_6.png}
\caption{正方二维布拉维格子 $p$。}
\end{figure}

我们已说过，三维布拉维格子是无穷的点阵列，其中每个点都有相同的环境与取向。\textit{晶胞}（unit cell）可以用多种方式定义，它由三个非共面向量给出，这三个向量构成平行六面体的棱。记号 $\mathbf{a}$、$\mathbf{b}$、$\mathbf{c}$ 用于表示这三个向量。通常总是这样选取晶胞矢量，使晶胞明显体现格子的对称性；这种惯用的晶胞取法可能不是初基的，也就是说它可能含有不止一个（等价的）格点。建立在布拉维格子三个基本矢量 $\mathbf{t}_{1}$、$\mathbf{t}_{2}$、$\mathbf{t}_{3}$（式 (1.5.1)）之上的平行六面体称为\textit{初基晶胞}（primitive unit cell）。在惯用晶胞与初基晶胞不一致的每个布拉维格子中，都可以建立两者之间的关系；例如《X 射线晶体学国际表》的表 2.2.2 即做此事。图 1.7 中画出的晶胞是惯用晶胞。14 个布拉维格子各自的 $\mathbf{t}_{1}$、$\mathbf{t}_{2}$、$\mathbf{t}_{3}$ 的具体规定将在第 3 章（表 3.1）给出。惯用晶胞与初基晶胞的区别可以用二维例子说明：二维初基长方布拉维格子 $p$ 的格点排在一系列长方形的角点上，惯用晶胞就是其中一个长方形（见图 1.8(\textit{a})），它同时也是该格子的一个初基晶胞。若在每个这样的晶胞中心加一个点，就得到新的布拉维格子——心行长方布拉维格子 $c$，见图 1.8(\textit{b})。图中画出体现长方对称性的惯用晶胞（虚线）和一个初基晶胞（实线）：惯用晶胞实际含两个格点，而初基晶胞只含一个。有时还使用另一种晶胞——\textit{维格纳--赛茨原胞}（Wigner--Seitz unit cell，Wigner and Seitz 1933）。取格子上任一点 $O$ 为原点，作垂直平分 $O$ 与其最近邻（有时还包括次近邻）连线的平面族，由这些平面围成的晶胞称为维格纳--赛茨原胞；这类原胞的图示可参见 Delaunay（1932）。

布拉维格子如我们所描述的那样，是只有位置而没有大小和形状的数学点的阵列。现在我们希望在讨论中纳入构成晶体的原子或分子。实际上，我们是在每个布拉维格点上安置一组相同的原子或分子，而这组原子或分子本身必须具有与该晶体布拉维格子点群对称性相容的对称性。晶体的一种对称操作将是某种复合操作：先作某个点群操作，再作布拉维格子的一个平移 $\mathbf{t}$（见式 (1.5.1)）。于是，一种特别简单的\textit{空间群}（space group）就这样产生：让布拉维格子的每一点都联系着同一晶系的一个点群。这样可以得到 $\sum_{i=1}^{7} n_{i}p_{i}$ 个空间群，其中 $n_{i}$ 是某晶系中点群的数目，$p_{i}$ 是该晶系中布拉维格子的数目，下标 $i$ 遍历七个晶系。例如，立方晶系有 3 种布拉维格子和 5 个立方点群，因此可以用这种方法导出 15 个立方空间群。这种形式的空间群称为\textit{点式}（symmorphic）空间群，它是平移群与一个点群的半直积。半直积群的理论将在第 4 章讨论。上式并不能立即涵盖所有点式空间群，原因有二。其一，可以把三方点群与六方布拉维格子相联系：因此对三方晶系必须取 $p_{i} = 2$，尽管三方晶系只有一种布拉维格子。其二，实际上有时一个点群在布拉维格子上可能存在两种不同且可区分的取向；这导致少数额外的点式空间群不规则地出现在表 1.7 的第 6 列中。点式空间群实际上共有 73 个。

在点式空间群中，点群操作与平移群操作基本可以分开，出现的点群操作本身就是晶体的对称操作。第二类空间群可以从点式空间群出发，把某些反射面换成滑移反射面、把某些转动轴换成螺旋轴而得到；它们通常分别称为滑移面与螺旋轴。滑移反射操作是一种复合操作：先关于平面 $m$ 作反射，再平移 $\mathbf{v}$——$\mathbf{v}$ 往往（但不必然）就在平面 $m$ 内。接连两次这样的滑移反射操作的结果，必然是布拉维格子平移群中的一个平移。类似地可以有你这样的 $n$ 阶螺旋轴：螺旋旋转操作是绕轴转过 $2\pi/n$ 后再平移 $\mathbf{v}$，$\mathbf{v}$ 往往（但不必然）沿转轴方向。若这个螺旋旋转操作接连施行 $n$ 次，结果必是布拉维格子平移群中的一个纯平移。加入滑移反射面与螺旋轴后，可以再导出 157 个空间群，称为\textit{非点式空间群}（non-symmorphic space groups）或\textit{非点式空间群}（asymmorphic space groups）。可以附加条件：每个滑移反射操作的平移部分必须位于反射面自身之内，每个螺旋旋转操作的平移部分必须沿转轴方向。若施加这一条件，则当滑移面或螺旋轴存在时，对称元素就不再都作用在晶体的同一点上。此时总可以把算符变换到同一点上作用，但平移将失去"沿螺旋轴方向"或"在滑移面内"的意义，可能变成某个奇怪的方向。前一种看法在直观上更令人满意，后一种看法则对系统的数学处理必不可少；我们因此采用如下约定：所有空间群操作都在一点上作用。157 个非点式空间群与 73 个点式空间群合计 230 个空间群。它们在各晶系中的分布见表 1.7。

\begin{table}[H]
\centering
\textbf{表 1.7}\quad \textit{230 个空间群的分类}

\smallskip
\begin{tabular}{|l|c|c|c|c|c|c|c|c|}
\hline
晶系 & $s$ & $n_{i}$ & $p_{i}$ & $n_{i}p_{i}$ & 第 6 列 & 第 7 列 & 第 8 列 & 第 9 列 \\ \hline
三斜 & 1 & 2 & 1 & 2 & & 2 & & 2 \\ \hline
单斜 & 2 & 3 & 2 & 6 & & 6 & 7 & 13 \\ \hline
正交 & -- & 3 & 4 & 12 & 1 & 13 & 46 & 59 \\ \hline
四方 & 4 & 7 & 2 & 14 & 2 & 16 & 52 & 68 \\ \hline
三方 & 3 & 5 & 2 & 10 & 3 & 13 & 12 & 25 \\ \hline
六方 & 6 & 7 & 1 & 7 & 1 & 8 & 19 & 27 \\ \hline
立方 & & 5 & 3 & 15 & & 15 & 21 & 36 \\ \hline
\textbf{总计} & -- & 32 & 15 & 66 & 7 & 73 & 157 & 230 \\ \hline
\end{tabular}

\smallskip
\parbox{0.92\textwidth}{\textit{表 1.7 的注.}
(i) $s$ 是主轴存在时的阶数。主轴上的操作可以是真转动，也可以是非真转动。
(ii) $n_{i}$ 是给定晶系的点群数目。注意 $\bar{6}$（$C_{3h}$）与 $\bar{6}2m$（$D_{3h}$）必须归入六方晶系。$p_{i}$ 是给定晶系的布拉维格子数目，三方晶系例外，取 $p_{i} = 2$：因为如正文所释，每个三方点群都可以与六方布拉维格子相联系，故该格子总计被计了两次。
(iii) 第 6 列是由于某些点群在布拉维格子上有可区分的不同定向而出现的点式空间群数目。第 7、8 列分别是给定晶系的点式与非点式空间群总数。第 9 列是给定晶系两类空间群的总数。}
\end{table}

\needspace{5\baselineskip}
\textit{空间群算符的 Seitz 记号}

表示一般空间群操作的一个方便记号是 $\{R \,|\, \mathbf{v}\}$，其中 $R$ 是点群算符，$\mathbf{v}$ 是与之相联系的平移向量；该算符的解释是：它作用在空间的点上，而这些点总是相对一组固定的轴给出，即它是\textit{主动}算符（见 1.4 节）。符号 $\{R \,|\, \mathbf{v}\}$ 由 Seitz（1936\textsuperscript{b}）引入，称为 \textit{Seitz 空间群符号}。由于这些算符是主动的，$\{R_{1} \,|\, \mathbf{v}_{1}\}$ 作用在向量 $\mathbf{r}$ 上产生 $\mathbf{r}'$：

\begin{equation*}
\{R_{1} \,|\, \mathbf{v}_{1}\}\mathbf{r} = \mathbf{r}' = R_{1}\mathbf{r} + \mathbf{v}_{1},
\tag{1.5.2}
\end{equation*}

其中 $R_{1}\mathbf{r}$ 是主动点群算符 $R_{1}$ 作用在 $\mathbf{r}$ 上得到的向量。这些算符的乘法有一条简单而重要的规则：

\begin{equation*}
\{R_{2} \,|\, \mathbf{v}_{2}\}\{R_{1} \,|\, \mathbf{v}_{1}\} = \{R_{2}R_{1} \,|\, \mathbf{v}_{2} + R_{2}\mathbf{v}_{1}\}
\tag{1.5.3}
\end{equation*}

把 $\{R_{2} \,|\, \mathbf{v}_{2}\}$ 作用在式 (1.5.2) 两边即可容易验证。由式 (1.5.3) 立即得到 $\{R_{1} \,|\, \mathbf{v}_{1}\}$ 的逆为 $\{R_{1}^{-1} | -R_{1}^{-1}\mathbf{v}_{1}\}$。全部 230 个空间群的 Seitz 算符已被多位作者整理（Faddeev 1964，Henry and Lonsdale 1965，Koptsik 1966，Kovalev 1965，Lyubarskii 1960，Miller and Love 1967），并列于第 3 章。

\textit{空间群标记的记号}

空间群首先按其所属晶系分类；进一步的分类取决于想传达多少信息。可以标记空间群所\textit{同形}（isogonal）的点群——即从空间群的操作中取出所有点群算符所组成的点群。\footnote{有些作者用 isomorphous（同晶）一词表示我们这里的 isogonal（同形）。} 这正是使用 Schönflies 记号（Hilton 1903）时的立场：空间群被赋予其同形点群的 Schönflies 标记，同形于同一有理点群的不同空间群则以一个任意的区分上标相区别——该上标仅仅反映 Schönflies（1891）推导空间群的先后次序。例如，与立方点群 $O_{h}$（$m3m$）同形的空间群有 10 个，分布在三种立方布拉维格子上：四个初基、四个面心、两个体心，记作 $O_{h}^{1}, O_{h}^{2}, \ldots, O_{h}^{10}$。这种记号有两个缺点。其一，一眼看不出空间群建立在哪种布拉维格子上：例如 $O_{h}^{8}$ 是面心的，$O_{h}^{6}$ 是体心的。其二，看不出空间群是点式还是非点式的。更随意的方案是给每个空间群一个 1 到 230 之间的编号，由 Astbury and Yardley（1924）提出，并被收入《X 射线晶体学国际表》（Henry and Lonsdale 1965）。

Schönflies 记号的这两个缺点在使用国际记号时都被克服了：国际记号不仅确定上述两点，还给出空间群中对称元素的性质与取向的某些信息（Hermann 1928\textsuperscript{a}，Mauguin 1931，Schiebold 1929）。做法是在国际点群符号（例如 $m3m$）前加一个描述布拉维格子的字母——$P$、$F$、$I$ 分别表示初基、面心、体心——于是得到 $Pm3m$、$Fm3m$、$Im3m$。这三个符号表示与点群 $m3m$（$O_{h}$）同形的三个点式空间群。当一个点群在布拉维格子上可以有多种取向时（例如四方晶系中的 $\bar{4}2m$（$D_{2d}$）），可以赋予国际空间群符号中字符次序以特定的含义；于是 $P\bar{4}2m$（$D_{2d}$）与 $P\bar{4}m2$（$D_{2d}$）是不同的空间群。对非点式空间群，符号按下述方式修改：凡被滑移面代替的反射面，符号中相应部分换成另一个表示滑移面的字母；例如 $Fd3m$（$O_{h}^{7}$）——金刚石结构的空间群——是面心的、非点式的，具有对角方向的滑移面。其他字母还有晶轴方向 $a$、$b$、$c$。对螺旋轴，空间群符号中表示该轴阶数的部分（例如 $P432$ 中的 4）换成同一符号加上一个表示螺旋轴平移量大小的后缀；于是 $P4_{1}32$（$O^{7}$）——与点群 432（$O$）同形的立方空间群之一——是初基的，其四次转动轴换成四次螺旋轴，即每次四次转动伴随晶格矢量在该方向四分之一长度的平移。国际记号并不完美，但它广为人知、经过仔细鉴定（Henry and Lonsdale 1965）并被广泛接受。这一记号的原始形式见旧版《X 射线晶体学国际表》（Bragg, von Laue, and Hermann 1935\textsuperscript{a}, \textsuperscript{b}），后来的修改见新版（Henry and Lonsdale 1965）；如何在鉴定空间群时使用这些表将在第 3 章（3.5 节）进一步讨论。本书提到点群或空间群时，一律用国际符号加括号内的 Schönflies 符号。

\textit{函数空间算符}

若 $R$ 是 Seitz 空间群算符的转动部分，即 $R$ 是一个点群操作，则每个算符 $R$ 通过如下类型的方程作用在空间的点 $\mathbf{r}$ 上（类似式 (1.4.8)）：

\begin{equation*}
R\mathbf{r} = \mathbf{r}'.
\tag{1.5.4}
\end{equation*}

若定义了标量函数 $f(\mathbf{r})$，则点 $\mathbf{r}$ 的移动会引起函数的相应变化。新函数满足：新函数在变换后点的取值等于原函数在原点的取值，即

\begin{equation*}
g(\mathbf{r}') = f(\mathbf{r}).
\tag{1.5.5}
\end{equation*}

习惯上记 $g = \bar{R}f$，即对每个算符 $R$ 定义了一个对应的\textit{函数空间算符}（function space operator）$\bar{R}$。由式 (1.5.4) 与 (1.5.5) 得

\begin{equation*}
\bar{R}f(\mathbf{r}) = f(R^{-1}\mathbf{r}).
\tag{1.5.6}
\end{equation*}

遗憾的是，多数作者并不为 $\bar{R}$ 使用单独的符号；事实上，只要仔细区分算符作用在函数上还是作用在向量上，就无须单独的符号。处理函数空间算符时必须按式 (1.5.6) 进行。下面的例子足以作为警示。

\textit{例} 1.5.1. 取 $\mathbf{r} = x\mathbf{i} + y\mathbf{j} + z\mathbf{k}$，定义 $f(\mathbf{r}) = x$，$g(\mathbf{r}) = y$，$h(\mathbf{r}) = z$，即向量 $\mathbf{r}$ 在三个坐标方向的投影。则

\begin{equation*}
\bar{C}_{4z}^{+}f(\mathbf{r}) = f(C_{4z}^{+}\mathbf{r}),
\tag{1.5.7}
\end{equation*}

它等于 $f(y\mathbf{i} - x\mathbf{j} + z\mathbf{k})$（见表 1.4 的注 (iii)）。自变量在 $\mathbf{i}$ 方向的投影是 $y$，因此

\begin{equation*}
\bar{C}_{4z}^{+}f(\mathbf{r}) = y = g(\mathbf{r}).
\tag{1.5.8}
\end{equation*}

于是在符号层面上既有 $C_{4z}^{+}\mathbf{i} = \mathbf{j}$，又有 $\bar{C}_{4z}^{+}f = g$。事实上，以基 $\langle \mathbf{i}, \mathbf{j}, \mathbf{k}|$ 表示算符 $C_{4z}^{+}$ 的矩阵，与以基 $\langle f, g, h|$ 表示函数空间算符 $\bar{C}_{4z}^{+}$ 的矩阵是相同的。在提醒过要小心之后，这倒是一个人们几乎不敢期望其为真的结果。

\textit{例} 1.5.2. 设 $\phi$ 是方位角（$x$ 轴与位置矢量在 $xy$ 平面上投影之间的夹角）。则

\begin{equation*}
C_{4z}^{+}\phi = \phi + \tfrac{1}{2}\pi
\tag{1.5.9}
\end{equation*}

但把 $\bar{C}_{4z}^{+}$ 作为函数空间算符时必须写

\begin{equation*}
\bar{C}_{4z}^{+}\cos\phi = \cos(C_{4z}^{-}\phi) = \cos(\phi - \tfrac{1}{2}\pi) = \sin\phi.
\tag{1.5.10}
\end{equation*}

式 (1.5.10) 也可以由式 (1.5.8) 推出：因为 $f(\mathbf{r}) = r\sin\theta\cos\phi$，$g(\mathbf{r}) = r\sin\theta\sin\phi$，而 $C_{4z}^{+}$ 不改变 $r$ 与 $\theta$。还应注意，由于 $\cos(\phi + \tfrac{1}{2}\pi) = -\sin\phi$，有 $C_{4z}^{+}\cos\phi \neq \cos(C_{4z}^{+}\phi)$。

\textit{主动与被动空间群算符}

如前所述，主动算符移动空间的点或位置矢量，所有矢量都相对一组固定的轴给出。本书使用的 Seitz 空间群算符是主动的，由式 (1.5.2) 定义：

\begin{equation*}
\{R_{a} \,|\, \mathbf{v}\}\mathbf{r} = R_{a}\mathbf{r} + \mathbf{v},
\tag{1.5.11}
\end{equation*}

其中用下标 $a$ 强调这是主动算符。也就是说，先作主动转动 $R_{a}$ 使 $\mathbf{r} \rightarrow R_{a}\mathbf{r}$，再对新位置矢量作平移 $+\mathbf{v}$，结果为 $(R_{a}\mathbf{r} + \mathbf{v})$。主动 Seitz 空间群算符的乘法规则已在式 (1.5.3) 中得到：

\begin{equation*}
\{S_{a} \,|\, \mathbf{w}\}\{R_{a} \,|\, \mathbf{v}\} = \{S_{a}R_{a} \,|\, \mathbf{w} + S_{a}\mathbf{v}\}.
\tag{1.5.12}
\end{equation*}

$\{R_{a} \,|\, \mathbf{v}\}$ 的逆为

\begin{equation*}
\{R_{a} \,|\, \mathbf{v}\}^{-1} = \{R_{a}^{-1} \,|\, -R_{a}^{-1}\mathbf{v}\}.
\tag{1.5.13}
\end{equation*}

被动算符移动空间的坐标轴，而空间的所有点从而所有矢量位置保持不动，于是每次操作之后矢量位置都相对一组新的坐标轴给出。也可以定义被动空间群算符（Johnston 1960）。为与主动 Seitz 算符区分，使用方括号 $[R_{p} \,|\, \mathbf{v}]$，对应于式 (1.5.11) 的定义规则是

\begin{equation*}
[R_{p} \,|\, \mathbf{v}]\mathbf{r} = R_{p}\mathbf{r} - R_{p}\mathbf{v}.
\tag{1.5.14}
\end{equation*}

也就是说，先把轴平移 $+\mathbf{v}$，使 $\mathbf{r} \rightarrow (\mathbf{r} - \mathbf{v})$，再转动新轴，最终结果为 $(R_{p}\mathbf{r} - R_{p}\mathbf{v})$，其中 $R_{p}$ 现在是被动转动。这些被动算符的乘法规则是

\begin{equation*}
[S_{p} \,|\, \mathbf{w}][R_{p} \,|\, \mathbf{v}] = [S_{p}R_{p} \,|\, \mathbf{v} + R_{p}^{-1}\mathbf{w}]
\tag{1.5.15}
\end{equation*}

而 $[R_{p} \,|\, \mathbf{v}]$ 的逆为

\begin{equation*}
[R_{p} \,|\, \mathbf{v}]^{-1} = [R_{p}^{-1} \,|\, -R_{p}\mathbf{v}].
\tag{1.5.16}
\end{equation*}

若 $X$ 是一个算符——无论主动还是被动——则与 $X$ 一起还有另一个算符 $\bar{X}$：$\bar{X}$ 是式 (1.5.4)--(1.5.6) 定义的函数空间算符。若 $f(\mathbf{r})$ 是位置的标量函数，则

\begin{equation*}
\bar{X}f(\mathbf{r}) = f(X^{-1}\mathbf{r}).
\tag{1.5.17}
\end{equation*}

必须认识到 $(\bar{X}f)$ 是一个函数的记号，正如 $f$ 也是函数的记号。式 (1.5.17) 并非随意规定，它陈述了如下事实：变换后的函数 $(\bar{X}f)$ 在变换后的点 $X\mathbf{r}$ 处的取值，等于 $f$ 在点 $\mathbf{r}$ 处的取值。Wigner（1959）强调过的另一点也必须再次强调：若 $Z = YX$，则要求 $\bar{Z}f = \bar{Y}\bar{X}f$，因为我们要在算符群与相应的函数空间算符群之间保持同构。证明只需反复应用式 (1.5.17)。记 $\bar{X}f = g$，则

\begin{align*}
\bar{Y}\bar{X}f(\mathbf{r}) &= \bar{Y}g(\mathbf{r}), \\
&= g(Y^{-1}\mathbf{r}) &&\text{（由式 (1.5.17)）}, \\
&= \bar{X}f(Y^{-1}\mathbf{r}) &&\text{（因 } \bar{X}f = g\text{）}, \\
&= f(X^{-1}Y^{-1}\mathbf{r}) &&\text{（由式 (1.5.17)）}, \\
&= f(Z^{-1}\mathbf{r}) &&\text{（因 } Z = YX\text{）}, \\
&= \bar{Z}f(\mathbf{r}) &&\text{（再次由式 (1.5.17)）}.
\end{align*}

于是如所要求的有 $\bar{Y}\bar{X}f = \bar{Z}f$。直接对式 (1.5.17) 作用 $\bar{Y}$ 是错误的，因为式 (1.5.17) 是关于两个函数在不同自变量取值下的数值关系，而算符不能作用于数（Altmann and Bradley 1965，Slater 1965\textsuperscript{a}）。

作为第一个例子，考虑被动空间群算符对函数 $\psi(\mathbf{r})$ 的作用。由式 (1.5.17) 得

\begin{align*}
\overline{[S_{p} \,|\, \mathbf{w}]}\,\psi(\mathbf{r}) &= \psi([S_{p} \,|\, \mathbf{w}]^{-1}\mathbf{r}) \\
&= \psi([S_{p}^{-1} \,|\, -S_{p}\mathbf{w}]\mathbf{r}) &&\text{（由式 (1.5.16)）} \tag{1.5.18}\\
&= \psi(S_{p}^{-1}\mathbf{r} + \mathbf{w}) &&\text{（由式 (1.5.14)）}. \tag{1.5.19}
\end{align*}

由函数空间算符的乘法规则与式 (1.5.19)，

\begin{align*}
\overline{[S_{p} \,|\, \mathbf{w}]}\,\overline{[R_{p} \,|\, \mathbf{v}]}\,\psi(\mathbf{r}) &= \overline{[R_{p} \,|\, \mathbf{v}]}\,\psi(S_{p}^{-1}\mathbf{r} + \mathbf{w}) \\
&= \psi(R_{p}^{-1}S_{p}^{-1}\mathbf{r} + R_{p}^{-1}\mathbf{w} + \mathbf{v}). \tag{1.5.20}
\end{align*}

从式 (1.5.19) 出发按如下方式“推导”式 (1.5.20) 是不正确的：

\begin{align*}
\overline{[S_{p} \,|\, \mathbf{w}]}\,\overline{[R_{p} \,|\, \mathbf{v}]}\,\psi(\mathbf{r}) &= \overline{[S_{p} \,|\, \mathbf{w}]}\,\psi(R_{p}^{-1}\mathbf{r} + \mathbf{v}) \tag{1.5.21}\\
&= \psi(R_{p}^{-1}(S_{p}^{-1}\mathbf{r} + \mathbf{w}) + \mathbf{v}) \tag{1.5.22}\\
&= \psi(R_{p}^{-1}S_{p}^{-1}\mathbf{r} + R_{p}^{-1}\mathbf{w} + \mathbf{v}). \tag{1.5.20}
\end{align*}

式 (1.5.21) 之所以错误，是因为如上所述 $[S_{p} \,|\, \mathbf{w}]$ 不能直接作用于数值关系 $[R_{p} \,|\, \mathbf{v}]\psi(\mathbf{r}) = \psi(R_{p}^{-1}\mathbf{r} + \mathbf{v})$；退一步说即使可以，式 (1.5.22) 也不对，并且直接违背数学惯例：函数空间算符作用于函数的整个自变量，而不是只作用于其中的 $\mathbf{r}$ 部分。

对主动空间群算符，与式 (1.5.19)、(1.5.20) 相应的方程是

\begin{equation*}
\overline{\{S_{a} \,|\, \mathbf{w}\}}\,\psi(\mathbf{r}) = \psi(S_{a}^{-1}(\mathbf{r} - \mathbf{w}))
\tag{1.5.23}
\end{equation*}

以及

\begin{equation*}
\overline{\{S_{a} \,|\, \mathbf{w}\}}\,\overline{\{R_{a} \,|\, \mathbf{v}\}}\,\psi(\mathbf{r}) = \psi(R_{a}^{-1}S_{a}^{-1}\mathbf{r} - R_{a}^{-1}S_{a}^{-1}\mathbf{w} - R_{a}^{-1}\mathbf{v}).
\tag{1.5.24}
\end{equation*}

式 (1.5.19) 与 (1.5.23) 加上 $S_{p} = S_{a}^{-1}$ 这一事实，蕴涵 $[S_{p} \,|\, \mathbf{w}] = \{S_{a} \,|\, \mathbf{w}\}^{-1}$——这正是我们应当期望的主动算符与被动算符之间的关系。

\textit{例} 1.5.3. 令 $\{E \,|\, \mathbf{t}\}$ 表示平移 $\mathbf{t}$。则

\begin{equation*}
\{E \,|\, \mathbf{t}\}\mathbf{r} = \mathbf{r} + \mathbf{t}.
\tag{1.5.25}
\end{equation*}

而若 $\{E \,|\, \mathbf{t}\}$ 作用在函数 $\psi(\mathbf{r})$ 而非向量 $\mathbf{r}$ 上，则有

\begin{equation*}
\overline{\{E \,|\, \mathbf{t}\}}\,\psi(\mathbf{r}) = \psi(\{E \,|\, \mathbf{t}\}^{-1}\mathbf{r}) = \psi(\mathbf{r} - \mathbf{t}).
\tag{1.5.26}
\end{equation*}

如前所述，本书一律使用主动点群算符与主动空间群算符；只是偶尔举例说明：若选用被动算符，我们的理论讨论需要作哪些修改。
